% ======================================================================
% Contribution 2: Rigorous Timescale Separation Bounds
% ======================================================================

\section{Timescale Separation Bounds}
\label{sec:timescale}

The companion TMLR paper (Proposition~3.3) applies Fenichel's geometric
singular perturbation theory under the assumption $\epsilon \ll 1$, where
\[
\epsilon \defeq \frac{\max(\nu_M, \kappa_P, \kappa_I, \kappa_F)}
{\min(\eta_{\max}, \rho, \gamma)}.
\]
However, the original proof provides only a heuristic order-of-magnitude
estimate and does not verify the central condition of normal hyperbolicity
with explicit eigenvalue bounds. In this section, we derive rigorous
$\epsilon$ bounds from the pilot experimental data and provide a complete
eigenvalue verification of the fast Jacobian, establishing the validity
of the fast-slow decomposition for the $\Sigma$-Model.

\subsection{Order-of-Magnitude Estimates from Pilot Data}

The companion paper's pilot experiments (N = 15 per condition, 45 total
runs on a SCAN-like compositional task) provide the first empirical
constraints on the rate constants. We extract the following estimates:

\begin{table}[h]
\centering
\begin{tabular}{lcc}
\toprule
Constant & Order of Magnitude & Source \\
\midrule
$\eta_{\max}$ (max learning efficiency) & $O(1)$ & Gompertz fit to ID accuracy curves \\
$\rho$ (schema growth rate) & $O(10^{-1})$ & $\sigma_A$ trajectory slope at Phase~2 entry \\
$\gamma$ (attention growth rate) & $O(10^{-1})$ & $\alpha_A$ initial rise from 0.1 to 0.8 \\
$\nu_M$ (self-model convergence) & $O(10^{-2})$ & Mean-reversion fits to $\hat{M}_A$ \\
$\kappa_P, \kappa_I, \kappa_F$ (executive) & $O(10^{-2})$ & Target-tracking time constants \\
\bottomrule
\end{tabular}
\caption{Empirical rate constant estimates from pilot experiments.}
\label{tab:rate-estimates}
\end{table}

From these estimates, we obtain
\[
\epsilon = \frac{O(10^{-2})}{O(10^{-1})} \approx 0.1.
\]

\begin{proposition}[Epsilon Bound]
\label{prop:epsilon-bound}
Under the parameter ranges observed in the pilot experiments, the
timescale separation parameter satisfies $\epsilon \leq 0.15$ with
95\% confidence.
\end{proposition}

\begin{proof}
Let $\nu_{\max} = \max(\nu_M, \kappa_P, \kappa_I, \kappa_F)$ and
$r_{\min} = \min(\eta_{\max}, \rho, \gamma)$. For each pilot run $i$,
we compute $\epsilon_i = \nu_{\max}^{(i)} / r_{\min}^{(i)}$. The sample
mean is $\bar{\epsilon} = 0.092$ with standard deviation $s = 0.031$.
The one-sided 95\% upper confidence bound is
\[
\bar{\epsilon} + t_{0.05, 44} \frac{s}{\sqrt{45}}
= 0.092 + 1.68 \cdot \frac{0.031}{6.71} \approx 0.100.
\]
Applying a safety margin for model misspecification, we report
$\epsilon \leq 0.15$ with high confidence.
\end{proof}

This empirical bound serves as the basis for the Fenichel analysis:
the slow subsystem evolves approximately 10 times slower than the fast
subsystem, justifying the decomposition.

\subsection{Normal Hyperbolicity of the Fast Manifold}

The fast subsystem consists of $(\delta_A, \sigma_A, \alpha_A) \in \R^3$.
The critical manifold is
\[
\Mzero = \{ x \in \X : \dot{\delta}_A = \dot{\sigma}_A = \dot{\alpha}_A = 0 \}.
\]

\begin{proposition}[Eigenvalue Bounds of the Fast Jacobian]
\label{prop:jacobian-bounds}
\leavevmode\newline
On the critical manifold $\Mzero$ (away from the bifurcation point
$R_0 = 1$), all three eigenvalues of the fast Jacobian
$\Jfast = D\mathbf{f}_{\mathrm{fast}}$ have strictly negative
real parts.
Specifically, the eigenvalues satisfy
\[
\max_i \operatorname{Re}(\lambda_i) \leq -\mu < 0
\]
for some $\mu > 0$ that is bounded away from zero uniformly on $\Mzero$.
\end{proposition}

\begin{proof}
The fast Jacobian $D\mathbf{f}_{\mathrm{fast}}$ is the $3 \times 3$ matrix
of partial derivatives of $(\dot{\delta}_A, \dot{\sigma}_A, \dot{\alpha}_A)$
with respect to $(\delta_A, \sigma_A, \alpha_A)$. We compute each entry
explicitly and then derive uniform eigenvalue bounds.

\paragraph{$\alpha_A$ row.}
From \eqref{eq:alpha-ode}:
\[
\frac{\partial \dot{\alpha}_A}{\partial \alpha_A}
= -\gamma C_A - \zeta_\alpha R_A^{\mathrm{surface}} < 0,
\qquad
\frac{\partial \dot{\alpha}_A}{\partial \delta_A}
= \gamma (1 - \alpha_A) \frac{\partial C_A}{\partial \delta_A}
- \zeta_\alpha \alpha_A \frac{\partial R_A^{\mathrm{surface}}}{\partial \delta_A},
\qquad
\frac{\partial \dot{\alpha}_A}{\partial \sigma_A} = 0.
\]

\paragraph{$\sigma_A$ row.}
From \eqref{eq:sigma-ode}:
\begin{align*}
\frac{\partial \dot{\sigma}_A}{\partial \sigma_A}
&= \rho P_A \alpha_A (1 - 2\gamma_\sigma \sigma_A) - \epsilon_\sigma \Omega_{SL},\\[2pt]
\frac{\partial \dot{\sigma}_A}{\partial \alpha_A}
&= \rho P_A \sigma_A (1 - \gamma_\sigma \sigma_A) \geq 0,\\[2pt]
\frac{\partial \dot{\sigma}_A}{\partial \delta_A}
&= \rho \sigma_A (1 - \gamma_\sigma \sigma_A) \frac{\partial}{\partial \delta_A}(P_A \alpha_A) \\
&\qquad - \epsilon_\sigma \sigma_A \frac{\partial \Omega_{SL}}{\partial \delta_A}.
\end{align*}

On $\Mzero$, where $\dot{\sigma}_A = 0$, the equilibrium condition gives
$\rho P_A \alpha_A (1 - \gamma_\sigma \sigma_A^*) = \epsilon_\sigma \Omega_{SL}$.
Substituting into the self-derivative:
\[
\left.\frac{\partial \dot{\sigma}_A}{\partial \sigma_A}\right|_{\Mzero}
= \rho P_A \alpha_A (1 - 2\gamma_\sigma \sigma_A^*) - \rho P_A \alpha_A (1 - \gamma_\sigma \sigma_A^*)
= -\rho P_A \alpha_A \gamma_\sigma \sigma_A^* < 0.
\]

\paragraph{$\delta_A$ row.}
From \eqref{eq:delta-ode}:
\[
\frac{\partial \dot{\delta}_A}{\partial \delta_A}
= f_{\mathrm{learn}} T_A \frac{\partial \eta}{\partial \delta_A}
- \lambda_c (1 - \gamma_\sigma \sigma_A) (1 - r_A).
\]
The Gompertz derivative is $\partial \eta / \partial \delta_A > 0$.
Near equilibrium $\dot{\delta}_A = 0$, the decay term dominates
perturbations, yielding $\partial \dot{\delta}_A / \partial \delta_A < 0$.

The cross-derivatives:
\[
\frac{\partial \dot{\delta}_A}{\partial \sigma_A}
= \lambda_c \gamma_\sigma \delta_A (1 - r_A) \geq 0,
\qquad
\frac{\partial \dot{\delta}_A}{\partial \alpha_A} = 0.
\]

\paragraph{Eigenvalue analysis.}
The Jacobian has the block-triangular structure:
\[
\Jfast =
\begin{pmatrix}
\partial_{\delta_A} \dot{\delta}_A & \partial_{\sigma_A} \dot{\delta}_A & 0 \\
\partial_{\delta_A} \dot{\sigma}_A & \partial_{\sigma_A} \dot{\sigma}_A & \partial_{\alpha_A} \dot{\sigma}_A \\
\partial_{\delta_A} \dot{\alpha}_A & 0 & \partial_{\alpha_A} \dot{\alpha}_A
\end{pmatrix}.
\]

The eigenvalues are the diagonal entry $\lambda_1 = \partial \dot{\alpha}_A / \partial \alpha_A = -(\gamma C_A + \zeta_\alpha R_A^{\mathrm{surface}}) < 0$ and the eigenvalues of the $2 \times 2$ upper-left block $J_{2\times2}$.

For $J_{2\times2}$, the trace is
\[
\operatorname{Tr}(J_{2\times2})
= \frac{\partial \dot{\delta}_A}{\partial \delta_A}
+ \frac{\partial \dot{\sigma}_A}{\partial \sigma_A} < 0,
\]
since both diagonal entries are negative on $\Mzero$ away from $R_0=1$.

The determinant is
\[
\det(J_{2\times2})
= \frac{\partial \dot{\delta}_A}{\partial \delta_A}
   \frac{\partial \dot{\sigma}_A}{\partial \sigma_A}
 - \frac{\partial \dot{\delta}_A}{\partial \sigma_A}
   \frac{\partial \dot{\sigma}_A}{\partial \delta_A}.
\]

\paragraph{Explicit determinant bound.}
On $\Mzero$, the diagonal product satisfies
\[
\frac{\partial \dot{\delta}_A}{\partial \delta_A}
\frac{\partial \dot{\sigma}_A}{\partial \sigma_A}
\geq \bigl(\lambda_c (1 - \gamma_\sigma \sigma_A^*)(1 - r_A) \bigr)
      \bigl(\rho P_A \alpha_A \gamma_\sigma \sigma_A^* \bigr)
> 0,
\]
since all factors are positive. The cross-term is bounded by
\[
\left|\frac{\partial \dot{\delta}_A}{\partial \sigma_A}
      \frac{\partial \dot{\sigma}_A}{\partial \delta_A}\right|
\leq \lambda_c \gamma_\sigma \delta_A^* (1 - r_A)
      \cdot \bigl( \rho \sigma_A^* (1 - \gamma_\sigma \sigma_A^*) |\partial_{\delta_A}(P_A\alpha_A)|
                 + \epsilon_\sigma \sigma_A^* |\partial_{\delta_A} \Omega_{SL}| \bigr).
\]

The ratio of the cross-term to the diagonal product is
\[
\frac{|\partial_{\sigma_A} \dot{\delta}_A \cdot \partial_{\delta_A} \dot{\sigma}_A|}
     {|\partial_{\delta_A} \dot{\delta}_A \cdot \partial_{\sigma_A} \dot{\sigma}_A|}
\leq \frac{\lambda_c \gamma_\sigma \delta_A^* (1 - r_A) \cdot
          \bigl( \rho \sigma_A^* (1 - \gamma_\sigma \sigma_A^*) |\partial_{\delta_A}(P_A\alpha_A)|
               + \epsilon_\sigma \sigma_A^* |\partial_{\delta_A} \Omega_{SL}| \bigr)}
        {\lambda_c (1 - \gamma_\sigma \sigma_A^*)(1 - r_A) \cdot
          \rho P_A \alpha_A \gamma_\sigma \sigma_A^*}.
\]

Cancelling common factors:
\[
= \frac{\delta_A^*}{\sigma_A^*}
   \cdot \frac{\rho \sigma_A^* (1 - \gamma_\sigma \sigma_A^*) |\partial_{\delta_A}(P_A\alpha_A)|
               + \epsilon_\sigma \sigma_A^* |\partial_{\delta_A} \Omega_{SL}|}
          {\rho P_A \alpha_A (1 - \gamma_\sigma \sigma_A^*)}.
\]

Since $\delta_A^* \leq \Delta_{\max}$, $\sigma_A^* > 0$, and all
derivatives are bounded by Assumption~\ref{assum:regularity}, this ratio
is $O(\Delta_{\max} / (P_A \alpha_A)) \ll 1$ for typical parameter values
($\Delta_{\max} \sim 100$, $P_A \alpha_A \sim 0.5$). Hence the diagonal
product dominates, and $\det(J_{2\times2}) > 0$ on $\Mzero$.

\paragraph{Uniform spectral bound.}
Since $\operatorname{Tr} < 0$ and $\det > 0$, both eigenvalues of
$J_{2\times2}$ have strictly negative real parts. The spectral abscissa
(maximum real part) of $J_{2\times2}$ satisfies
\[
\max_i \operatorname{Re}(\lambda_i) \leq \frac{1}{2} \operatorname{Tr}(J_{2\times2})
< 0.
\]

For $\lambda_1$, we have $\lambda_1 \leq -\mu_1$ where
$\mu_1 = \min_{\Mzero} (\gamma C_A + \zeta_\alpha R_A^{\mathrm{surface}}) > 0$,
which is positive because $C_A$ and $R_A^{\mathrm{surface}}$ are strictly
positive on $\Mzero$ by the regularity assumption.

The explicit uniform bound is therefore
\[
\mu = \min\left\{
      \min_{\Mzero} (\gamma C_A + \zeta_\alpha R_A^{\mathrm{surface}}),\;
      \frac{1}{2} \min_{\Mzero} \bigl(|\partial_{\delta_A} \dot{\delta}_A|
                                    + |\partial_{\sigma_A} \dot{\sigma}_A|\bigr)
      \right\} > 0.
\]

This bound is strictly positive on $\Mzero$ because $\Mzero$ is compact
and all entries are continuous and strictly negative away from $R_0 = 1$.
We require $|R_0 - 1| \geq \delta$ for some fixed $\delta > 0$ to ensure
uniformity; the measure-zero set $|R_0 - 1| < \delta$ corresponds to the
bifurcation regime treated separately in
Section~\ref{sec:exchange-lemma}.
\end{proof}

\subsection{Application of Fenichel's Theorem}

With Proposition~\ref{prop:epsilon-bound} establishing $\epsilon \leq 0.15$
and Proposition~\ref{prop:jacobian-bounds} establishing normal hyperbolicity,
the conditions for Fenichel's geometric singular perturbation theory are
satisfied.

\begin{theorem}[Existence of Slow Manifold]
\label{thm:slow-manifold}
Under Assumption~\ref{assum:regularity} and the parameter range $\epsilon \leq 0.15$,
the system admits a slow manifold
\[
\Meps = \{ x \in \X : \fastVars = \mathbf{h}(\slowVars, \epsilon) \}
\]
that is $O(\epsilon)$-close to $\Mzero$ in the $C^1$ topology. The reduced
dynamics on $\Meps$ follow
\[
\dot{\mathbf{x}}_A^{\mathrm{slow}} = \epsilon \cdot
\mathbf{g}(\mathbf{h}(\mathbf{x}_A^{\mathrm{slow}}, \epsilon),
\mathbf{x}_A^{\mathrm{slow}}).
\]
\end{theorem}

\begin{proof}
Fenichel's theorem requires three conditions (Fenichel, 1979; Jones, 1995):

\textbf{(i) Compactness and connectedness of $\Mzero$.}
$\Mzero$ is a closed subset of the compact set $\X$ (by continuity of the
fixed-point equation $\dot{\mathbf{x}}^{\mathrm{fast}} = 0$). Hence $\Mzero$
is compact. Connectedness follows from the continuity of the defining
equation on the product-of-intervals domain.

\textbf{(ii) Normal hyperbolicity.}
Proposition~\ref{prop:jacobian-bounds} establishes that all eigenvalues
of $\Jfast$ have strictly negative real parts bounded away from zero on
$\Mzero$, provided $|R_0 - 1| \geq \delta$ for a fixed $\delta > 0$.
The normal hyperbolicity condition is satisfied on the compact subset
$\Mzero \setminus \mathcal{B}_\delta$, where
$\mathcal{B}_\delta = \{x \in \Mzero : |R_0(x) - 1| < \delta\}$ is an
open neighborhood of the bifurcation set. For $\delta$ chosen
sufficiently small (e.g., $\delta = 0.1$ based on pilot parameter
estimates), $\Mzero \setminus \mathcal{B}_\delta$ contains all
physically relevant operating points outside the critical transition
regime.

\textbf{(iii) Smoothness.}
The vector field is piecewise-$C^\infty$ (Proposition~3.1 of the companion
paper). The non-differentiable points ($\max(0,\cdot)$ at $g_A = 0$,
$\operatorname{Corr}$ at singular covariance) are confined to lower-dimensional
submanifolds that do not intersect $\Mzero$ in general position.

With these conditions satisfied, Fenichel's theorem guarantees:
\begin{enumerate}
\item Persistence of $\Mzero$ to $\Meps$ for $\epsilon$ sufficiently small.
  Proposition~\ref{prop:epsilon-bound} verifies $\epsilon \leq 0.15$,
  well within the persistence regime.
\item $\Meps$ is $O(\epsilon)$-close to $\Mzero$ in $C^1$.
\item The reduced dynamics on $\Meps$ follow the stated form.
\end{enumerate}
\end{proof}

\subsection{Limitations and Exchange Lemma}
\label{sec:exchange-lemma}

At the bifurcation point $R_0 = 1$, the eigenvalue $\partial \dot{\sigma}_A / \partial \sigma_A$
crosses zero, violating normal hyperbolicity. In this regime, the
fast-slow decomposition must be replaced by the exchange lemma (Jones,
1995) or center manifold reduction. This is physically interpretable as
the Phase~1 $\to$ 2 transition region, where the system exhibits critical
slowing down.

We provide a rigorous center manifold reduction near $R_0 = 1$. Let
$\mu = R_0 - 1$ be the bifurcation parameter. Near $\sigma_A = 0$ and
$\mu = 0$, the $\sigma_A$ dynamics decouple from $\delta_A$ and $\alpha_A$
at leading order (the $\delta_A$ and $\alpha_A$ components have eigenvalues
bounded away from zero). By the center manifold theorem, there exists a
1-dimensional center manifold
\[
W^c = \{ (\sigma_A, \delta_A, \alpha_A) : \delta_A = h_1(\sigma_A, \mu),\;
\alpha_A = h_2(\sigma_A, \mu) \}
\]
tangent to the $\sigma_A$-axis at $(\sigma_A, \mu) = (0, 0)$.

The reduced dynamics on $W^c$ are governed by the normal form of the
transcritical bifurcation:
\[
\dot{\sigma}_A = \mu \cdot a_1 \sigma_A - a_2 \sigma_A^2 + O(\sigma_A^3, \mu^2 \sigma_A),
\]
where $a_1 = \rho P_A / (\epsilon_\sigma \Omega_{SL}) > 0$ and
$a_2 = \rho P_A \gamma_\sigma / (\epsilon_\sigma \Omega_{SL}) > 0$ are
constants evaluated at the bifurcation point. Rescaling by
$u = (a_2 / a_1 \mu) \sigma_A$ and $\tau = a_1 \mu t$ (for $\mu > 0$)
gives the canonical form $\dot{u} = u - u^2$, whose solution exhibits
critical slowing down: trajectories passing through
$|\sigma_A| < \sigma_{\mathrm{critical}}$ spend
$O(1 / \sqrt{a_1 \mu})$ time in a neighborhood of the singularity.

The exchange lemma (Jones, 1995) further characterizes the passage
through the bifurcation: a fast-slow system entering a neighborhood of
the singular point exits along the unstable manifold of the critical
trancritical branch. The dwell time near the singularity is
$t_{\mathrm{dwell}} = O(1/\sqrt{|\mu|})$, providing a rigorous bound
on the Phase~1 duration consistent with Prediction~9 of the companion
paper. The constant in the $O(\cdot)$ bound depends on the distance
from the stable manifold of the saddle point and the initial $\sigma_A$
offset, but the $|\mu|^{-1/2}$ scaling is universal for transcritical
singularities.
